% !PRINTED none
% Part-title / new-Abschnitt opening page: "Siebenundzwanzigster Abschnitt. Algebraische und
% Abelsche Differentiale." followed directly (same page) by the heading "§ 197. Differentiale
% in Ω." and the opening paragraphs and formulas of that section. No running-head page number
% is visible on this page in the scan.
\begin{center}
\gesperrt{Siebenundzwanzigster Abschnitt.}

\textbf{Algebraische und Abelsche Differentiale.}

\rule{0.3\linewidth}{0.4pt}

\textbf{\S\,197. \gesperrt{Differentiale in $\Omega$.}}
\end{center}

Sind $z$ und $z_1$ zwei Variable in $\Omega$ von den Ordnungen $n$, $n_1$ mit den Unterecken $\mathfrak{U}$, $\mathfrak{U}_1$, mit den Verzweigungsecken: $\mathfrak{Z}$, $\mathfrak{Z}_1$, den Verzweigungszahlen $w$, $w_1$, so ist nach \S\ 195, (12)
\[
\dfrac{dz}{dz_1} = \dfrac{\mathfrak{Z}\mathfrak{U}_1^2}{\mathfrak{Z}_1\mathfrak{U}^2}.
\]

Setzt man
\weq{1}{\omega = \dfrac{\mathfrak{A}\mathfrak{U}^2}{\mathfrak{B}\mathfrak{Z}}, \quad \omega_1 = \dfrac{\mathfrak{A}\mathfrak{U}_1^2}{\mathfrak{B}\mathfrak{Z}_1},}
so wird hiernach
\weq{2}{\omega\,dz = \omega_1\,dz_1.}

Ist $\mathfrak{A}\mathfrak{U}^2$ \"aquivalent mit $\mathfrak{B}\mathfrak{Z}$, so ist $\omega$ eine Funktion in $\Omega$, und $\omega_1$ ist gleichfalls eine Funktion in $\Omega$. Sind $a$, $b$ die Ordnungen von $\mathfrak{A}$ und $\mathfrak{B}$, so ist [\S\ 196, (2)]
\weq{3}{\begin{aligned}
b+w &= a+2n,\\
a &= b+2p-2,
\end{aligned}}
wenn $p$ das Geschlecht des K\"orpers $\Omega$ ist.

Wir setzen jetzt in einer neuen symbolischen Bezeichnung
\weq{4}{dz = \dfrac{\mathfrak{Z}}{\mathfrak{U}^2},}
\weq{5}{\omega\,dz = dJ = \dfrac{\mathfrak{A}}{\mathfrak{B}},}
und nennen diese Ausdr\"ucke die zum K\"orper $\Omega$ geh\"origen \gesperrt{Differentiale}.

Im Z\"ahler und Nenner $\mathfrak{A}$, $\mathfrak{B}$ eines Differentials k\"onnen gemeinschaftliche Faktoren zugef\"ugt oder weggelassen werden. Haben $\mathfrak{A}$ und $\mathfrak{B}$ keinen gemeinschaftlichen Teiler, so hei\ss{}t $\mathfrak{A}$ das \gesperrt{Obereck}, $\mathfrak{B}$ das \gesperrt{Untereck} des Differentials $dJ$.
